\begin{proof}
We express the Ricci tensor of the Riemannian metric in terms of the Lee form, its covariant derivative and its divergence, and then compare two integral identities for this derivative, one from the contracted Bianchi identity and one from the Bochner formula. Norms, traces and contractions are taken with respect to $g$. For real tangent vectors $U,V$, put
\[
 A=\lambda^\sharp,\qquad B=JA,\qquad
 S=\nabla\lambda,\qquad
 S^+(U,V)=\tfrac12\bigl(S(U,V)+S(JU,JV)\bigr).
\]
The tensor $S$ is symmetric because $\mathrm{d}\lambda=0$, and $S^+$ is its orthogonal projection onto the $J$-invariant symmetric tensors. Thus $S^-=S-S^+$ and $|S|^2=|S^+|^2+|S^-|^2$. Put $\delta=\mathrm{d}^*\lambda$, so that $\delta=-\operatorname{tr}_gS=-\operatorname{div}A$, and write $J\lambda=-\lambda\circ J$, so that $\lambda^{\mathrm c}=J\lambda/(4\pi)$ and $(J\lambda)^\sharp=B$.

\emph{Step 1: the Ricci tensor.} We first show that
\begin{equation}\label{eq:lck-ricci-real}
 \operatorname{Ric}_g
 =nS^+-(n-1)S+\tfrac12\bigl(\delta-|\lambda|^2\bigr)g
       +\tfrac12\lambda\otimes\lambda
       +\tfrac n2 J\lambda\otimes J\lambda.
\end{equation}
To prove this identity, take a local function $f$ with $\lambda=\mathrm{d}f$ and put $\widetilde g=\mathrm{e}^{-f}g=\mathrm{e}^{2\varphi}g$, where $\varphi=-f/2$. The formula for the Ricci tensor under a conformal change, in real dimension $2n$, is
\[
 \operatorname{Ric}_{\widetilde g}
 =\operatorname{Ric}_g-(2n-2)\bigl(\nabla\mathrm{d}\varphi-\mathrm{d}\varphi\otimes\mathrm{d}\varphi\bigr)
       -\bigl(\operatorname{tr}_g\nabla\mathrm{d}\varphi+(2n-2)|\mathrm{d}\varphi|^2\bigr)g.
\]
Since $\nabla\mathrm{d}\varphi=-S/2$ and $\operatorname{tr}_g\nabla\mathrm{d}\varphi=\delta/2$, this gives
\[
 \operatorname{Ric}_{\widetilde g}
 =\operatorname{Ric}_g+(n-1)S
       +\tfrac{n-1}{2}\lambda\otimes\lambda
       -\tfrac12\bigl(\delta+(n-1)|\lambda|^2\bigr)g.
\]
The local metric $\widetilde g$ is K\"ahler, and its volume form is $\mathrm{e}^{-nf}$ times that of $g$. Since the Chern--Ricci form of $g$ vanishes, the Ricci form of $\widetilde g$, written with the unnormalized operator $\mathrm i\partial\bar\partial$, is
$n\mathrm i\partial\bar\partial f=\tfrac n2\mathrm{d}(J\lambda)$.

Applying \eqref{eq:lck-complex-connection} to
$(\nabla_UJ\lambda)(V)=-S(U,JV)-\lambda((\nabla_UJ)V)$ and antisymmetrizing yields
\[
 \mathrm{d}(J\lambda)(U,JV)
 =2S^+(U,V)-|\lambda|^2g(U,V)
             +\lambda(U)\lambda(V)+J\lambda(U)J\lambda(V).
\]
For a K\"ahler metric, the Ricci tensor is obtained from the Ricci form by evaluation on $(U,JV)$. Comparing the last two displays with the conformal formula proves \eqref{eq:lck-ricci-real}. The calculation is local, but the resulting tensor identity holds globally. Since $\operatorname{tr}_gS^+=\operatorname{tr}_gS=-\delta$, taking traces in \eqref{eq:lck-ricci-real} gives the scalar curvature
\[
 \operatorname{scal}_g=(n-1)\delta-\tfrac{n-1}{2}|\lambda|^2.
\]

\emph{Step 2: the Bianchi identity.} All integrals below are taken with respect to the Riemannian volume of $g$, which we omit from the notation. Since
$S(A,A)=\tfrac12A(|\lambda|^2)$ and $\operatorname{div}A=-\delta$,
\begin{equation}\label{eq:lck-aa-zero}
 \int_X S(A,A)=\frac12\int_X|\lambda|^2\delta.
\end{equation}
By the contracted Bianchi identity, the divergence of the Ricci tensor is half the differential of the scalar curvature. Integrating by parts twice therefore gives
\[
 \int_X\langle\operatorname{Ric}_g,S\rangle
 =-\frac12\int_XA(\operatorname{scal}_g)
 =-\frac12\int_X\operatorname{scal}_g\,\delta.
\]
We substitute \eqref{eq:lck-ricci-real} and the scalar curvature into this identity. Since $S^+$ is the orthogonal projection of $S$ onto the $J$-invariant symmetric tensors and $\langle g,S\rangle=-\delta$, and since \eqref{eq:lck-aa-zero} expresses $\int_X|\lambda|^2\delta$ through $\int_XS(A,A)$, we obtain
\begin{equation}\label{eq:lck-bianchi-integral}
 n\int_X|S^+|^2-(n-1)\int_X|S|^2+\frac{n-2}{2}\int_X\delta^2
 +\frac{4-n}{2}\int_XS(A,A)+\frac n2\int_XS(B,B)=0.
\end{equation}

\emph{Step 3: the Bochner formula.} Pairing the Weitzenb\"ock formula for one-forms with $\lambda$ and integrating gives, since $\mathrm{d}\lambda=0$,
\[
 \int_X\delta^2=\int_X|S|^2+\int_X\operatorname{Ric}_g(A,A).
\]
When \eqref{eq:lck-ricci-real} is evaluated on $(A,A)$, its two quartic terms cancel and the term $J\lambda\otimes J\lambda$ vanishes. Thus
\[
 \operatorname{Ric}_g(A,A)
 =(1-\tfrac n2)S(A,A)+\tfrac n2S(B,B)+\tfrac12|\lambda|^2\delta.
\]
Together with \eqref{eq:lck-aa-zero}, this gives
\begin{equation}\label{eq:lck-bochner-integral}
 \int_X\delta^2=\int_X|S|^2+\frac{4-n}{2}\int_XS(A,A)+\frac n2\int_XS(B,B).
\end{equation}
The last two terms are the same in \eqref{eq:lck-bianchi-integral} and \eqref{eq:lck-bochner-integral}. Eliminating them leaves
\[
 n\int_X|S^+|^2-n\int_X|S|^2+\frac n2\int_X\delta^2=0.
\]
Since $|S|^2-|S^+|^2=|S^-|^2$, this is the asserted identity.
\end{proof}
